\section{Introduction and independent problem formulation}
\label{sec:branch_introduction}
A time-indexed family of rolling-origin forecast-validation loss
functions can have multiple local minima, and the locations of
these minima can vary as new observations become available.
This study asks whether historical minima can be connected
across origins and converted into a useful forecast-time
smoothness decision. The question is distinct from merely
minimizing a pooled forecast-validation loss.

For an observed fitting window $x_t\in\mathbb R^L$,
difference order $1\le d<L$ and forecast horizon $h$,
define $D_d\in\mathbb R^{(L-d)\times L}$ and
$Q=D_d^\top D_d$. For $0\le\lambda<\infty$,
$H_\lambda=(I+\lambda Q)^{-1}$ is a well-defined
symmetric positive definite smoother; at the limiting
polynomial projection it is defined by continuity.
Use the attained-smoothness index $S\in[0,1]$ and a
fixed finite-difference polynomial continuation matrix
$G_{d,h}$ to construct forecasts. If $z_t$ denotes
the subsequently observed $h$-step block, write
\begin{equation}
\label{eq:branch_origin_mse}
F_t(S)=\frac{1}{h}
\|z_t-G_{d,h}H(S)x_t\|_2^2.
\end{equation}
All $z_t$ used at operational origin $T$ must satisfy
$t+h\le T$; $z_T$ is not observed at selection time.

An ordinary historical benchmark selects
\begin{equation}
\label{eq:F_S}
\widehat S_T^{\mathrm{pool}}
\in\operatorname*{arg\,min}_{S\in[0,1]}
\frac{1}{M}\sum_{m=1}^M F_{t_m}(S).
\end{equation}
The dynamic object instead retains each surface's
locally optimal smoothness values and tests whether
their temporal trajectories contain predictive
information that this pooled reduction ignores.
The resulting branch policies require separate
assumptions, outer evaluation, and evidence; they
do not inherit forecast optimality from the pooled
criterion.

In all procedures, a selected $S$ is used to refit
$H(S)y_{T-L+1:T}$ at the current origin, before applying
$G_{d,h}$. No historical trend estimates are averaged
or carried into the final forecast. The retained
historical observations alone determine the decision.
